\begin{proof}[Proof of \cref{obj:lem:pilot-radius-second-moment}]
Since \(n\geq1\) and \(d\geq16\), we have \(m>0\) and \(L>0\). For each \(\zeta\in\mathcal Z\), the localization intervals in \cref{obj:def:jf-estimator} give
\[
c_{\zeta j}=\frac{N'_{\zeta j}}m,\qquad
h_{\zeta j}
=H_0\left(\sqrt{\frac{N'_{\zeta j}L}{m^2}}+\frac Lm\right),\qquad
R_{\zeta j}
=\frac{c_{\zeta j}+h_{\zeta j}
-\max\{0,c_{\zeta j}-h_{\zeta j}\}}2.
\]
Both \(c_{\zeta j}\) and \(h_{\zeta j}\) are nonnegative. The maximum in the last expression lies between \(c_{\zeta j}-h_{\zeta j}\) and \(c_{\zeta j}+h_{\zeta j}\); hence \(0\leq R_{\zeta j}\leq h_{\zeta j}\). Applying \((x+y)^2\leq2(x^2+y^2)\) yields
\[
R_{\zeta j}^2
\leq 2H_0^2\left\{
\frac{N'_{\zeta j}L}{m^2}
+\left(\frac Lm\right)^2
\right\}.
\]
% lean: CausalSmith.Stat.DiscreteBudgetvalueCurve.pilotRadius_sq_le_count_envelope

There are four categories in \(\mathcal Z\). Cauchy–Schwarz and the coordinate bound therefore give, pointwise,
\[
S_j^2
=\left(\sum_{\zeta\in\mathcal Z}R_{\zeta j}\right)^2
\leq4\sum_{\zeta\in\mathcal Z}R_{\zeta j}^2
\leq8H_0^2\left\{
\frac L{m^2}\sum_{\zeta\in\mathcal Z}N'_{\zeta j}
+4\left(\frac Lm\right)^2
\right\}.
\]
% lean: CausalSmith.Stat.DiscreteBudgetvalueCurve.pilotRadius_sum_sq_integral_le_cellScale

Each pilot count is integrable under its Poisson law, with
\[
\mathbb E\!\left[\sum_{\zeta\in\mathcal Z}N'_{\zeta j}\right]
=\sum_{\zeta\in\mathcal Z}m q_{\zeta j}
=mp_j.
\]
The pointwise bound thus proves integrability of \(S_j^2\) and gives
\[
\mathbb E[S_j^2]
\leq8H_0^2\left\{
\frac{p_jL}{m}+4\left(\frac Lm\right)^2
\right\}.
\]
Finally, \(p_jL/m\geq0\) and \(L/m\geq0\), so
\[
\frac{p_jL}{m}+4\left(\frac Lm\right)^2
\leq4\left(\sqrt{\frac{p_jL}{m}}+\frac Lm\right)^2
=4v_j^2.
\]
Combining the two inequalities proves \(\mathbb E[S_j^2]\leq32H_0^2v_j^2\).
% lean: CausalSmith.Stat.DiscreteBudgetvalueCurve.pilotRadius_sum_sq_integral_le_cellScale
\end{proof}
